\section{问题重述与分析}

微网利用储能转移电量，并根据外网收费规则决定提前采购及补购。本题的目标是在满足负载、容量及功率限制的条件下尽可能节省实际购电费用。问题一给定单日输入，建立确定性日前联合线性规划作为物理基础\cite{kong2016,silva2020}。问题二引入负载与光伏的不确定性，五倍紧急电价和无退款的计划费用构成非对称预测误差成本，采用历史残差场景进行采购\cite{abunima2022,marino2018,grover2018}。问题三在规定发布时点允许付费修订，进一步区分常规购电承诺与10分钟储能控制\cite{palma2013,silvente2015,li2016,jani2022}。问题四保留上述结构，通过联合残差刻画价格与净负荷的关联\cite{li2020}。

四问的主路线为“确定性LP、非对称结算下的场景采购、交易约束下的双时间尺度控制、价荷联合不确定下的滚动调度”。本文用收费机制解释采购行为，并核对优化参考计划和实际控制所产生的费用。CVaR、备用及内层LP均接受实际账单检验。风险项与备用在经济主方案中关闭；一月校准四案均选中因果贪心执行器，因此内层LP作为已实现的候选及对照，不能写成必然优于原控制规则的主结论。

\section{模型假设与符号说明}

\subsection{基本假设}

\begin{enumerate}
 \item 将源文件功率解释为以标记时刻为右端点的10分钟平均功率，以$\Delta t=1/6$小时换算电量。这是时间离散的建模解释，不是赛题明确给出的事实。忽略线路损耗及自放电，充电、放电效率分别为$\eta_c=\eta_d=0.9$，往返效率为0.81。
 \item 不售电，富余电量允许免费弃置。主口径下计划购电全额付费，取消部分不退款；退款解释仅用于敏感性分析。不引入题目未给出的折旧、碳成本或投资费用。
 \item 储电量保持在1200--10800 kWh，充放电功率均不超过5000 kW。问题一首末储电量均为6000 kWh；其他问题仅在1月1日初始化，以后连续传递状态。
 \item 每10分钟可观测当期负载、光伏、电价及储电量并完成平衡。紧急购电可补足剩余缺口，不另设题目未给出的容量上限。未来实际值不可提前使用，附件4全天实际价格不作为日前已知输入。
\end{enumerate}

\subsection{主要符号}

购电量和电池外侧充放电量统一使用kWh，以$p$表示电价、$c$表示充电量。
\begin{table}[htbp]
\centering\small
\caption{主要符号及含义}\label{tab:symbols}
\begin{tabularx}{\linewidth}{p{0.18\linewidth}Xl}
\toprule
符号 & 含义 & 单位\\
\midrule
$D,t,s$ & 日期、日内时段$t=0,\ldots,143$、场景索引 & ---\\
$L_{D,t},R_{D,t},n_{D,t}$ & 实际负载功率、光伏功率、净负荷电量 & kW、kW、kWh\\
$p_{D,t}$ & 实际交易电价 & 元/kWh\\
$g_{D,t},a_{D,t}$ & 零时原计划、最终有效常规购电量 & kWh\\
$c_t,d_t,e_t$ & 电池外侧充电量、放电量、时段初储电量 & kWh\\
$h_t,w_t$ & 紧急购电量、弃置电量 & kWh\\
$\pi_s,Q_s$ & 场景概率、场景剩余时域紧急费用 & ---、元\\
$\rho_t,\lambda$ & 外层参考备用量、CVaR权重 & kWh、---\\
$\beta_T,k_\rho$ & 终端偏离惩罚系数、备用缩放系数 & 元/kWh、---\\
$E_{\min},E_{\max},U$ & 储电量下上限、单段电量上限$5000/6$ & kWh\\
\bottomrule
\end{tabularx}
\end{table}

\section{数据处理与因果预测}

\subsection{统一物理时间索引与原模板映射}

每天保留144段，定义
\begin{equation}
 I_t=[t/6,(t+1)/6),\qquad n_{D,t}=(L_{D,t}-R_{D,t})/6.
 \label{eq:time_energy}
\end{equation}
源标记00:10对应物理00:00--00:10，$0{:}00{+}1$对应物理23:50--24:00，最后一段归入当天结算。附件5按要求保留原有表头，从0:10--0:20开始，以次日0:00--0:10结束；按源样本列序写入，不因表头再平移数据。

映射冻结于\texttt{data/time\_mapping.json}：每个$t$同时记录源时间、物理区间、模板标签、单元格位置及四小时分组。指定时段表、储能汇总、紧急区间和工作簿均从同一索引生成。例如$t=60$对应源10:10、物理10:00--10:10、模板10:10--10:20；题目指定的10:00--10:10结果取$t=60$，而不按模板文字反查。紧急区间按物理连续时段合并，四小时汇总按连续24段求和。这一区分使论文自然日结算与原模板显示均可追溯。

\subsection{一月费用校准与固定运行前缀}

修订模型以1月15--31日为验证窗口，按实际结算费用选择预测器及控制器，2月1日--12月31日334天用于评价与填报。所有同一问题的验证候选从已冻结旧版策略在1月15日的实际储电量出发，各自连续运行至1月31日。候选的期末电量同时报告，但不作为虚拟退款冲减验证费用，因此这是账单比较，并非库存价值已经中和的比较。

候选集合预先限定：问题二、三的负载在前一周、最近2个及4个同星期日均值间选择；问题二的历史光伏在3、7、14日均值间选择。动态电价两案继承对应固定价问题选定的负载及光伏预测器，再在上述三类同星期方法中选择价格预测器。采用逐项坐标搜索，先固定28场景、无风险和无备用的内层LP配置选择预测，再比较贪心与内层LP控制器；滚动两案另比较启用跨日续接的LP。该有限候选搜索不等同于预测与控制全部组合的全局最优。风险及备用关闭、终端系数倍率1等未搜索设置明确固定。四案控制器均由一月账单选中贪心规则，结果及库存边界见下表。

\begin{table}[htbp]\centering\small
\caption{一月验证冻结的预测配置与执行结果}
\label{tab:v2_forecast_selection}
\begin{tabularx}{\linewidth}{LLLLRR}\toprule
问题 & 负荷 & 光伏 & 电价 & 验证费/元 & 末库存/kWh\\\midrule
问题2 & 同星期2日 & 最近3日 & 已知电价 & 917318.59 & 7155.31\\
问题3 & 同星期2日 & 附件3发布值 & 已知电价 & 919340.82 & 7143.30\\
问题4-2 & 同星期2日 & 最近3日 & 同星期4日 & 1042234.20 & 7298.41\\
问题4-3 & 同星期2日 & 附件3发布值 & 同星期2日 & 1043044.13 & 7308.72\\
\bottomrule\end{tabularx}
\par\smallskip\parbox{\linewidth}{\footnotesize 验证期为1月15日至31日；所有主案选择规则执行、28个历史场景，额外CVaR与备用均关闭。预测器先以统一储能LP进行逐坐标费用比较，再比较执行器，属于有限候选搜索。}
\end{table}

正式2月初态承接各问题已经冻结的旧版1月运行末态，同一问题所有修订策略使用相同初态；不会用1月底才选定的方法回写1月运行以改变2月库存。随后实际末态逐日传递。旧版1月记录仅作为共同运行前缀，1月验证模拟用于选型，两者用途不同。尽管2--12月数据不进入此次选型目标，修订研究已参考过旧版该时段结果，因此本文称其为冻结规则后的顺序回测，不称为从未接触的独立测试。

\subsection{预测更新及误差评价}

令$X$表示负载或价格，同星期预测使用已完成的历史日：
\begin{equation}
 \widehat X_{D,t}=\frac{1}{|\mathcal W_D|}\sum_{j\in\mathcal W_D}X_{j,t},\quad
 \mathcal W_D=\{D-7i:i=1,\ldots,K\}\cap\{1,\ldots,D-1\}.
 \label{eq:weekday_forecast}
\end{equation}
移动均值则取此前$K$个完整日。历史不足使用已有样本，无相应星期日时退回近7日均值，首日固定使用附件1先验。历史预测的每个日期均只读取此前数据；全年数组的预先生成不改变这一信息限制。

滚动问题使用当次已发布的24点光伏预报，将发布时刻$r$后第1至24个小时作为预报横坐标，以发布时刻已知光伏值为锚点线性插值到10分钟右端点。日内负载与动态电价用此前2小时的平均预测偏差$b_r$修正：
\begin{equation}
 \widehat X_{r,k}=\max\{\ell_X,\widehat X^{(0)}_{r,k}
 +\operatorname{clip}(b_r e^{-k/36},-0.3\widehat X^{(0)}_{r,k},0.3\widehat X^{(0)}_{r,k})\}.
 \label{eq:issue_correction}
\end{equation}
其中$k$为发布后的步数，负载下界$\ell_X=0$，预测价格下界为$10^{-6}$元/kWh。6小时衰减尺度和30\%幅度事先固定。只在0、6、12、18时更新这组预测，其他时刻的内层LP候选仅替换当前一步实际值，不使用尚未发布的光伏预报。

预测误差报告$\mathrm{RMSE}=\sqrt{N^{-1}\sum_i(X_i-\widehat X_i)^2}$及$\mathrm{MAE}=N^{-1}\sum_i|X_i-\widehat X_i|$。表中2--12月指标只作描述，不重新选型；滚动问题另外报告零时发布光伏预报的误差，避免与未用于其采购的历史光伏预测混淆。
\begin{table}[htbp]\centering\small
\caption{已选预测器的误差指标（光伏发布值取零时版本）}
\label{tab:v2_forecast_metrics}
\begin{tabularx}{\linewidth}{LLRRRR}\toprule
问题 & 变量 & 一月MAE & 一月RMSE & 评价期MAE & 评价期RMSE\\\midrule
问题2 & 负荷/kW & 122.51 & 152.09 & 171.91 & 240.54\\
问题2 & 光伏/kW & 100.38 & 214.94 & 154.46 & 313.48\\
问题3 & 负荷/kW & 122.51 & 152.09 & 171.91 & 240.54\\
问题3 & 光伏/kW & 135.37 & 266.89 & 197.49 & 383.03\\
问题4-2 & 负荷/kW & 122.51 & 152.09 & 171.91 & 240.54\\
问题4-2 & 光伏/kW & 100.38 & 214.94 & 154.46 & 313.48\\
问题4-2 & 电价/(元/kWh) & 0.03390 & 0.04364 & 0.04886 & 0.06963\\
问题4-3 & 负荷/kW & 122.51 & 152.09 & 171.91 & 240.54\\
问题4-3 & 光伏/kW & 135.37 & 266.89 & 197.49 & 383.03\\
问题4-3 & 电价/(元/kWh) & 0.03434 & 0.04386 & 0.04515 & 0.06383\\
\bottomrule\end{tabularx}
\par\smallskip\parbox{\linewidth}{\footnotesize 一月指1月15日至31日；评价期指2月至12月。评价期误差仅作结果描述，不参与选型；价格列与功率列单位不同。}
\end{table}
